\documentclass{article}
\usepackage[margin=1in]{geometry}
\usepackage{mathtools, amsfonts, amsthm}

\usepackage{listings}


\title{HW 0}
\author{Sam Ly}

\begin{document}
\maketitle

\section*{Chapter 1}

\begin{enumerate}
    \item {
        Markdown, Jinja2, Makefile, SASS, CSS, HTML.
    }

    \item {
        \lstinputlisting[language=Java]{./ch1/java/Main.java}
    }

    \item {
        \lstinputlisting[language=C]{./ch1/c/main.c}
    }

\end{enumerate}

\section*{Chapter 3}

\begin{enumerate}
    \item {
        One potential edge case that wasn't defined in the book is nested classes.
        For example, if you were to make a Linked List class, the Node class may 
        be an inner class that doesn't need to exist outside of the Linked List class 
        scope. 

        Attempting to use inner classes actually leads to a compile error:

        \begin{lstlisting}
class LinkedListStack {
    class Node {
        init(data, next) {
            self.data = data;
            self.next = next;
        }
    }

    init() {
        self.head = nil;
    }
    // ...
}

---------------------------------------------
[line 2] Error at 'class': Expect method name.

        \end{lstlisting}

        The correct way to implement this is to write:
        \begin{lstlisting}
class Node {
    init(data, next) {
        this.data = data;
        this.next = next;
    }
}

class LinkedListStack {

    init() {
        this.head = nil;
    }

    push(data) {
        if (this.head == nil) {
            this.head = Node(data, nil);
        }
        else {
            var next = Node(data, this.head);
            this.head = next;
        }
    }

    pop() {
        if (this.head == nil) {
            return nil;
        } 
        else {
            var out = this.head;
            this.head = this.head.next;
            return out;
        }
    }
}
        \end{lstlisting}

        This behaves as I expect since Lox is meant to be a simple language.
        Nested classes introduce a lot of complexity into the implementation.
    }

    \item {
        One unanswered question I have relates to the idea of ``truthy-ness''.
        In languages like Python and JavaScript, certain can evaluate to values 
        when used as the predicate for an ``if'' statement. Also, these two 
        languages even have different implementations of truthy-ness!

        Example in JS:
        \lstinputlisting[]{./truthyness.js}

        Example in Python:
        \lstinputlisting[language=Python]{./truthyness.py}
    }

    \item {
        The most important feature that is missing in my opinion is the ability 
        to define return types for functions. I am a big fan of type safety, 
        but I am aware that type safety is quite difficult to implement
    }
\end{enumerate}



\end{document}